\emph{(i).} In the flat scale the Schouten tensor vanishes, so the tractor connection of Fact \claimref{dark-energy.imported-frozen-point-geometry} reads $\nabla_a(\sigma, \mu_b, \rho) = (\partial_a\sigma - \mu_a,\ \partial_a\mu_b + \eta_{ab}\rho,\ \partial_a\rho)$. A tractor is parallel exactly when $\rho$ is a constant $r$, $\mu_a = \partial_a\sigma$ and $\partial_a\partial_b\sigma = -r\eta_{ab}$, whose solutions are $\sigma = s + \eta(m, x) - \tfrac{r}{2}Q(x)$ with constants $s, m$; the space of parallel tractors is six-dimensional and the connection is flat. The norm is $\eta^{-1}(d\sigma, d\sigma) + 2\sigma r$; with $d\sigma = m^\flat - r x^\flat$ this is $Q(m) - 2r\eta(m,x) + r^2Q(x) + 2rs + 2r\eta(m,x) - r^2Q(x) = 2rs + Q(m)$, constant. By the Fact, a parallel tractor of norm $\Lambda/3 \neq 0$ is the scale tractor of an Einstein scale with that $\Lambda$ on the region where $\sigma \neq 0$, the metric being $\sigma^{-2}\eta$. The check \texttt{check\_parallel\_frozen\_points\_and\_norms} verifies, with exact multivariate polynomials in $x$, parallelism and the constant norm for four tractors, three of norm $2$ and one of norm $5$, and that the tractor form has signature $(2,4)$.

\emph{(ii).} If $\sigma_A = c\,\sigma_B$ for a constant $c$ the two quadratic polynomials would be proportional; the check \texttt{check\_two\_frozen\_points\_give\_two\_metrics} verifies that their coefficient vectors have rank two for the three equal-norm pairs among $I_A = (1,0,1)$, $I_B = (1,(1,0,0,0),\tfrac12)$, $I_C = (2,(0,1,0,0),\tfrac34)$, so the conformal factor between $\sigma_A^{-2}\eta$ and $\sigma_B^{-2}\eta$ is not constant (it takes different values at three rational points near $0$, where both scales are nonzero). Both metrics are Einstein with $\Lambda = 6$ by (i), and both lie in the conformal class of $\eta$. The pairing $\ip{I_A}{I_B}$ is constant because both tractors are parallel and the connection preserves the tractor form; at $x = 0$ it is $\eta(m_A, m_B) + s_A r_B + r_A s_B = 0 + \tfrac12 + 1 = \tfrac32$, which is not $\pm 2$, so $I_B \neq \pm I_A$.

\emph{(iii).} The Lie algebra of $SO(2,4)$ is the solution space of $A^{T}G + GA = 0$ for the $6 \times 6$ Gram matrix $G$ of the tractor form at a point, of dimension $15$; adding $A v = 0$ for $v = I_A$ cuts it to dimension $10$ (the stabiliser of a non-null vector is $\mathfrak{so}$ of its five-dimensional complement), and adding $A v = 0$ for $v = I_B$ as well cuts it to $6$, $\mathfrak{so}$ of the four-dimensional complement of the plane spanned by $I_A, I_B$, which is nondegenerate since $Q(I_A)\,Q(I_B) - \ip{I_A}{I_B}^2 = 4 - \tfrac94 \neq 0$. The check \texttt{check\_pair\_has\_one\_gauge\_invariant} computes these three dimensions by exact rank. Each frozen point of fixed norm ranges over a five-dimensional quadric, so pairs form a space of dimension $10$; the orbit of the pair has dimension $15 - 6 = 9$, leaving one invariant (the complementary result that where the Weyl tensor has no kernel vector every parallel tractor is a multiple of one is filed alongside, synthesis.frozen-point-unique-where-weyl-has-no-kernel), and $\ip{I_A}{I_B}$ is invariant and not constant on the space of pairs, so it is that invariant. Proposition \claimref{dark-energy.point-infinity-has-no-physical-fluctuation} concerns one pair $(I, V)$ subject to the Stelle-West constraints, whose fluctuations the nine broken generators exhaust; with a second frozen point the same count leaves the pairing as a physical datum, which that proposition does not address.

\emph{(iv).} For each tower the functional $Q(I_r)^{-1}\int\epsilon(I_r, V_r, F, F)$ is quadratic in $F$. Its first variation under $\delta I_r$, $\delta V_r$ (including the variation of the prefactor) is a sum of terms containing $F \otimes F$, and under $\delta\nabla$, with $\delta F = D\,\delta\nabla$, it is $2Q(I_r)^{-1}\int \epsilon(I_r, V_r, D\,\delta\nabla, F)$, linear in $F$. Every term vanishes where $F = 0$, so on $U$ the configuration with any frozen points $I_r$ and any $V_r$ obeying the constraints is stationary for the sum over towers.

\emph{Conclusion.} On $U$ every premise used in (i)-(iv) holds for the two-tower configuration $(I_A, I_B)$, which violates the metric clause of Premise \claimref{dark-matter.p1-single-frozen-point}; so that clause does not follow from them there. Since the parallel tractors on a region of vanishing Weyl part form a six-dimensional space, the pairing takes a continuum of values, and P1 is the further choice $I_B = \pm I_A$, which fixes it at $\pm\Lambda/3$.
